\section{Proofs}


\subsection*{A.1\quad Preliminaries and notation}
We work in the Cantor topology on $\Aomega$, whose basic open sets are the
cylinders $u\Aomega$ for $u\in\Astar$. A set $P\subseteq\Aomega$ is \emph{closed}
iff it contains all its limit points, where $x$ is a limit point of $P$ if every
prefix $x_{\le n}$ is a prefix of some element of $P$. We use repeatedly the
classical characterization of safety.

\begin{lemma}[Safety is topological closure]\label{lem:safety-closed}
$P\subseteq\Aomega$ is a safety property (Def.~\ref{def:safety}) if and only if
$P$ is closed. Equivalently, $x\in P$ iff every finite prefix $x_{\le n}$ is a
\emph{good prefix}, i.e., $x_{\le n}\notin\Bad(P)\Astar$.
\end{lemma}
\begin{proof}
If $P$ is safety and $x\notin P$, then $x$ has a bad prefix $x_{\le n}$, so the
open cylinder $x_{\le n}\Aomega$ is disjoint from $P$; hence the complement of $P$
is open and $P$ is closed. Conversely, if $P$ is closed and $x\notin P$, then $x$
is not a limit point, so some prefix $x_{\le n}$ extends into $P$ nowhere, i.e.,
$x_{\le n}\Aomega\cap P=\emptyset$; that prefix is a bad prefix, so $P$ is safety.
The prefix characterization is the contrapositive: $x\in P$ iff no prefix is bad.
\end{proof}

Recall (Def.~\ref{def:gate}) that a gate $g$ decides from committed effect
history $\tau_{<t}\in(\A\setminus\{\noop\})^{*}$, proposal $a_t$, and
$\Pi(\tau_{<t},a_t)$ (in the oracle case $g$ is computable relative to $\Pi$, so it may also use the bits $\Pi$ returned at earlier committed steps; for the fixed $\Pi$ its decision is still a function of $\tau_{<t}$ and $a_t$). We assume $\Pi$ reads this same effect history (equivalently,
an interface defined on full committed traces must be $\noop$-blind).
We keep $\Pi$ explicit: it is an available oracle interface,
not a promise that its bits are definable in $\Th$. A
\emph{$(\Th;\Pi)$-gate} runs a deterministic keyed-counter register automaton (Def.~\ref{def:safety}), whose
guards combine $\Th$ tests with queries to these bits, once and online over the committed effect history, and decides each
proposal by the one transition the automaton takes on it from the configuration reached (labeled $\allow$, $\block$, or
$\subst(t)$ for a term $t$ over the proposal's fields, the registers, and constants). A
\emph{proposal stream} is a map from committed histories to proposed actions;
together with $g$ it determines the committed trace. We write $\noop$ for the
neutral no-op of Def.~\ref{def:actions}. Because policies are scored on \emph{effect}
traces (Def.~\ref{def:actions}), the recognizer $\mathcal R$ of $\Good(P)$ reads only
effect actions: committing $\noop$ produces no effect and leaves $\mathcal R$'s state
unchanged (so $\delta(q,\noop)=q$), and a trace is safe for $P$ iff no finite prefix of
its effect trace is bad and an infinite effect trace lies in $P$ (for safety $P$ the
second clause follows from the first by Lemma~\ref{lem:safety-closed}). Table~\ref{tab:notation} collects the notation.

\begin{table}[t]
\caption{Notation.}
\label{tab:notation}
\centering\footnotesize
\begin{tabularx}{\columnwidth}{@{}>{\raggedright\arraybackslash}p{0.30\columnwidth}>{\raggedright\arraybackslash}X@{}}
\toprule
Symbol & Meaning\\
\midrule
$\A,\Astar,\Aomega$ & action alphabet; finite and infinite traces\\
$\noop$, $\erase$ & neutral no-op action; external-effect projection (deletes $\noop$)\\
$P,\Bad(P),\Good(P)$ & policy; minimal bad prefixes; good prefixes\\
$\Th$, $\Pi$ & guard theory; fixed vector of oracle predicate bits\\
$\RA[\Th;\Pi]$ & languages recognized by keyed-counter register automata with $\Th$-guards and oracle bits $\Pi$\\
$\pi_i,\hat\pi_i^{\theta}$ & ground-truth predicate; thresholded judge\\
$\alpha(\theta),\beta(\theta)$ & judge false-alarm and miss rates\\
$\ROC(\beta)$, $\ROC_I(\beta)$ & ROC frontier: least $\alpha$ at miss $\beta$ (given $I$)\\
$g,\Dec$ & gate; decisions $\{\allow,\block,\subst\}$\\
$R,U$ & episode unsafe rate; utility loss\\
$F_H(\delta)$ & controlled episode frontier (Def.~\ref{def:controlled})\\
$h_t$, $\psi$, $\mu$ & gate's observable history; gate policy; logging policy (Def.~\ref{def:controlled}, Prop.~\ref{prop:explore})\\
$F_0,F_1$, $Q_\rho$ & compliant and violating score laws; attacked violation-score law\\
$o$ & untrusted context the judge reads (Def.~\ref{def:adversary})\\
$r(\theta)$, $\hat\theta$ & per-step conditional miss risk; calibrated threshold\\
$H,\delta$ & horizon; target episode risk\\
$\eta(\rho;\theta),\bar\eta(\rho)$ & robustness margin at $\theta$; uniform margin $\sup_\theta\eta(\rho;\theta)$\\
\bottomrule
\end{tabularx}
\end{table}

\subsection*{A.2\quad T1: oracle characterization}
\noindent\textbf{Theorem~\ref{thm:t1}} (Oracle characterization)\textbf{.}
\emph{Fix the available oracle interface $\Pi$ and a nonempty policy $P$. Then $P$ is precisely
enforceable by a deterministic $(\Th;\Pi)$-gate iff it is safety and
$\Good(P)\in\RA[\Th;\Pi]$; the witness simulates that recognizer and refuses
exactly the one-step extensions it rejects.}

\begin{proof}
\emph{($\Leftarrow$).} Let $\mathcal R$ be a deterministic keyed-counter register automaton
(Def.~\ref{def:safety}) recognizing $\Good(P)$ with guards in $(\Th;\Pi)$; below, $q$ ranges over its
configurations (control state, registers, key maps, counters), $q_0$ is the initial one, $\delta$ the
deterministic update, and $B$ the absorbing bad configurations. Define the gate $G$ that maintains
$q_t=\delta^{*}(q_0,c_{\le t})$ and, on
proposed $a$ at state $q$, commits $a$ (moving to $\delta(q,a)$) if
$\delta(q,a)\notin B$, and otherwise commits $\noop$ (a \emph{block}, leaving the
state at $q$).

\emph{Soundness.} Blocking commits $\noop$ (no effect; the recognizer does not move);
an allowed action moves within $\Good(P)$ by construction, since $G$ blocks exactly the
one-step extensions entering $B$, and $q_0\notin B$ because $\varepsilon\in\Good(P)$ for nonempty $P$ (the only use of nonemptiness: for $P=\emptyset$, $\Good(P)=\emptyset$ is recognizable yet no gate is sound). So the committed run visits only good states and every
effect prefix $\erase(c_{\le t})\in\Good(P)=\Astar\setminus\Bad(P)\Astar$ is not bad.
Hence the committed trace is safe for $P$ (Def.~\ref{def:actions}): if the effect trace is
infinite it lies in $P$ by Lemma~\ref{lem:safety-closed}; if the gate eventually blocks
forever, the finite effect trace still has no bad prefix. This holds for every proposal
stream.

\emph{Transparency.} Suppose the proposed $\omega$-trace $w$ lies in $P$. By
Lemma~\ref{lem:safety-closed} every prefix of $w$ is good, so running $\mathcal R$
on $w$ never reaches $B$; hence $G$ never blocks and commits $w$ unchanged.

\emph{Complexity.} Each step evaluates the enabled guarded transition on the
current data, registers, and supplied oracle bits. This is polynomial when
guard evaluation and updates are polynomial; no satisfiability solver is
implicitly assumed.

\emph{($\Rightarrow$).} We work in the effect alphabet $\A\setminus\{\noop\}$, over which
$P$ and its traces live. Suppose the $(\Th;\Pi)$-gate $g$ precisely enforces $P$ but $P$ is not
safety. By Lemma~\ref{lem:safety-closed}, $P$ is not closed, so it has a limit point
$x\notin P$ that is a doomed \emph{effect} trace: for every $n$ there is an effect trace
$z^{(n)}\in P$ with $z^{(n)}_{\le n}=x_{\le n}$. Let $S$ be the \emph{open-loop} stream
proposing $x_1,x_2,\dots$ in turn (a fixed sequence, independent of the gate's actions). We
show by induction on $t$ that $g$ allows $x_t$, so the committed effect trace is $x$.

Base case $t=0$ is immediate. For the step, assume $g$ has allowed $x_1,\dots,x_{t-1}$, so the
committed effect history is $x_{<t}$ and the proposal is $x_t$. Consider the open-loop stream
$S^{(t)}$ proposing $z^{(t)}_1,z^{(t)}_2,\dots$; its proposed effect trace is $z^{(t)}\in P$,
so by \emph{prefixwise} transparency (Def.~\ref{def:gate}) the committed effect prefix equals
$z^{(t)}_{\le k}$ at every step $k$. Thus $g$ blocks or alters nothing on $S^{(t)}$, and in
particular allows $z^{(t)}_t=x_t$ from the committed effect history $z^{(t)}_{<t}=x_{<t}$.
Because $g$ conditions on the \emph{effect} history (Def.~\ref{def:gate}), its decision here is
the value $g(x_{<t},x_t,\Pi(x_{<t},x_t))$, identical under $S$ and $S^{(t)}$; being
deterministic, $g$ also allows $x_t$ under $S$. By induction $g$ commits the infinite
effect trace $x\notin P$, contradicting soundness (clause (ii) of Def.~\ref{def:actions}).
Note that $x$ has no bad prefix (every $x_{\le n}$ extends into $P$), which is why soundness
must constrain infinite effect traces and not only their finite prefixes. Hence $P$ is safety.

It remains to recognize $\Good(P)$. Simulate the finite control, registers, and
committed-history updates of $g$ while reading a proposed finite word $u$, and
move to a rejecting sink at the first step whose committed action differs from the proposal (an identity $\subst(a)$ counts as $\allow$). This is a
deterministic $(\Th;\Pi)$ register automaton. If $u\in\Good(P)$, then $u$ has a
completion in $P$ (otherwise $u$ itself would be bad), and prefixwise transparency
forces every symbol of $u$ to be allowed. Conversely, if $u$ contains a first bad
prefix $va$, transparency on a compliant completion of $v$ forces $g$ to allow
$v$, while soundness forbids allowing $a$ from that history; thus the simulation
rejects. Its language is exactly $\Good(P)$, proving
$\Good(P)\in\RA[\Th;\Pi]$.
\end{proof}

\begin{remark}[General computable gates]
Dropping the register restriction, the same limit-point argument shows an
arbitrary computable gate with interface $\Pi$ precisely enforces a nonempty
$P$ iff $P$ is safety and good-prefix membership is decidable relative to $\Pi$.
$(\Th;\Pi)$-gates are the register-recognizable subclass.

For a randomized gate, quantify separately: for every (possibly adaptive)
proposal strategy, soundness holds with probability one over gate randomness;
and for every fixed compliant proposal trace, prefixwise transparency holds
with probability one. On each compliant $z^{(t)}$ agreeing with a doomed limit
trace $x$ through $t$, transparency makes the conditional decision at history
$x_{<t}$ equal to $\allow$ almost surely. The countable intersection of these
probability-one events makes the gate commit all of $x$ almost surely,
contradicting almost-sure soundness. Thus randomization does not extend the
safety boundary (although the register-representability converse above is stated
only for deterministic gates).
\end{remark}

\medskip
\noindent\textbf{Proposition~\ref{prop:strict}} (A strict hierarchy)\textbf{.}
\emph{(i)~The gate-enforceable policies are a strict subclass of the edit-automaton
enforceable policies (edit automata given the same oracle bits $\Pi$). (ii)~If $\Pi$ is action-local and recipients range
over an infinite equality-only data sort, then no $(\Th;\Pi)$-gate without key maps (whatever its registers, global counters,
and counter guards and updates) enforces the per-recipient cap of one, which a $(\Th_0;\Pi)$-gate with one key-local flag
enforces.}

\begin{proof}
Every gate-enforceable policy is safety with a good-prefix language decidable relative to $\Pi$
(Theorem~\ref{thm:t1}, and the remark above for general computable gates); an edit automaton given the same oracle bits $\Pi$ enforces each such policy by suppressing any action
that would enter $\Bad(P)$ and re-emitting nothing (safety properties are a subclass of
the renewal properties edit automata enforce \citep{ligatti2005edit,ligatti2009nonsafety}),
giving containment. For strictness,
let $\A$ contain irreversible $\mathsf{pay}$ and $\mathsf{confirm}$ and let
\[
  P^{\ast}=\{w\in(\A\setminus\{\noop\})^{\omega}:\text{every }\mathsf{pay}\text{ has a later }\mathsf{confirm}\}.
\]
\emph{$P^{\ast}$ is not safety.} The effect trace $\mathsf{pay}^{\omega}$ violates
$P^{\ast}$, yet no finite prefix is bad: any prefix $\mathsf{pay}^{k}$ extends to
$\mathsf{pay}^{k}\,\mathsf{confirm}^{\omega}\in P^{\ast}$.
By Lemma~\ref{lem:safety-closed}, $P^{\ast}$ is not safety, so by Theorem~\ref{thm:t1} no
gate enforces it. \emph{Self-contained impossibility.} Each $\mathsf{pay}$ must be
decided irrevocably. On the open-loop stream $\mathsf{pay}^{\omega}$, suppose the gate
has allowed the first $t-1$ payments; the committed effect history $\mathsf{pay}^{t-1}$
and the proposal $\mathsf{pay}$ are then exactly those of the compliant stream
$\mathsf{pay}^{t}\,\mathsf{confirm}^{\omega}\in P^{\ast}$, on which prefixwise
transparency forces $\allow$. So the gate allows every $\mathsf{pay}$ and commits the
infinite effect trace $\mathsf{pay}^{\omega}\notin P^{\ast}$; the payments are
irreversible, no later decision can repair the violation, and the gate is unsound. Substitution does not help:
replacing $\mathsf{pay}$ by a semantically distinct action fails transparency,
and replacing it by an equivalent $\mathsf{pay}$ returns to the first case.
\emph{Edit-enforceable.} An $\omega$-policy counts as edit-enforceable when some extension of it to finite executions is. Containment above judges finite executions as good prefixes; that extension of $P^{\ast}$ (which has no bad prefix) accepts every finite execution and is not renewal, since $\mathsf{pay}^{\omega}\notin P^{\ast}$ has infinitely many valid prefixes. If instead a finite execution is judged valid iff each of its $\mathsf{pay}$s has a later $\mathsf{confirm}$ within it, $P^{\ast}$ is a renewal property; an edit automaton
suppresses each unconfirmed $\mathsf{pay}$, and every action after it, into a buffer and emits the buffer in order upon the next $\mathsf{confirm}$
(and never if $\mathsf{confirm}$ does not occur), which is edit-enforceable by the
renewal-property theorem of \citet{ligatti2009nonsafety}. Thus $P^{\ast}$ separates
the classes. The comparison is at the level of enforceable policy \emph{classes}, each in
its own model's semantics: edit automata enforce $P^{\ast}$ by buffering, which presumes
actions can be delayed and re-emitted; irreversibility removes exactly that power, which is
why the classical class is unattainable for a pre-execution gate over committed actions,
not a defect of the gate.

\emph{Part (ii).} Let $P_1$ be the set of infinite effect traces in which no two $\mathsf{send}$ actions name the same
recipient (other actions are unconstrained). It is safety, its bad prefixes being exactly the finite effect traces with two
$\mathsf{send}$s to one recipient, and it is nonempty because the recipient sort is infinite. A $(\Th_0;\Pi)$-gate enforces $P_1$
by Theorem~\ref{thm:t1}: a (K) component maps each recipient to a flag set by its first send and rejects a send whose flag is
set; it is trim, since a send to an unflagged recipient is always available, so its non-bad language is exactly $\Good(P_1)$.
Now let $g$ be a $(\Th;\Pi)$-gate without key maps, with $r$ registers and finitely many global counters, over an
equality-only recipient sort: a recipient enters its guards and updates only through equality with registers, with the other
fields of the current action, and with the finitely many constants of the automaton (guard constants, initial register
contents, and constants in substitution terms), and through the oracle bits, which for action-local $\Pi$ depend on the
proposed action alone. Fix the non-recipient fields $\bar f$ of $\mathsf{send}$. As $\Pi$ takes finitely many values, some
infinite set $D'$ of recipients, avoiding the automaton's constants and the data values in $\bar f$, gives sends with equal
bits. Feed $\mathsf{send}(d_1,\bar f),\ldots,\mathsf{send}(d_{r+1},\bar f)$ for pairwise distinct $d_i\in D'$; this prefix extends
to the compliant stream that continues with fresh recipients, so prefixwise transparency allows every send. The configuration
reached stores at most $r$ data values, so some $d_j$ is in no register. Let $e\in D'$ be neither among the $d_i$ nor in any
register. In this configuration, $\mathsf{send}(d_j,\bar f)$ and $\mathsf{send}(e,\bar f)$ satisfy exactly the same guards
(neither recipient equals a register, a constant, or a value in $\bar f$, and their bits agree), so $g$ takes the same
transition on both and outputs the same decision term. Transparency on the compliant stream that continues with
$\mathsf{send}(e,\bar f)$ and further fresh recipients forces that term to commit $\mathsf{send}(e,\bar f)$ unchanged. As $e$
equals no register value and no constant, the term is $\allow$ or a substitution that copies the proposed recipient, and on
$\mathsf{send}(d_j,\bar f)$ either one commits $\mathsf{send}(d_j,\bar f)$, completing a bad prefix and contradicting soundness.
The argument is the pumping argument of \citet{kaminski1994finite}. It needs determinism and the register-automaton form: a
general computable gate (Remark~\ref{rem:general-gates}) enforces $P_1$ without key maps, and so does a register-free gate
given the history-dependent bit $\mathbf{1}[\tau\cdot a\in\Good(P_1)]$.
\end{proof}

\medskip
\noindent\refstepcounter{proposition}\textbf{Proposition~\theproposition} (Substitution: same class; a utility gain only with a judge)\textbf{.}
\emph{(i)~Adding $\subst(a')$ does not change the class in T1, and in the oracle regime a
transparent gate never alters a compliant proposal, so substitution cannot lower $U$ there.
(ii)~Consider a judge-driven gate (relaxing the block-only clause of
Assumption~\ref{ass:benign}) that, whenever its judge flags a proposal at good prefix $u$, can
compute from its own inputs a rewrite with $ua'\in\Good(P)$ (certified without the
ground-truth bits) and conditional loss at most $\ell_{\rm sub}\in[0,1]$ relative to unit loss
for a block. Then substituting instead of blocking reduces the false-block loss at good prefixes from
$\alpha(\theta)$ to at most $\ell_{\rm sub}\alpha(\theta)$ at unchanged miss.}
\begin{proof}
(i)~Run the T1 recognizer. When $ua$ would be rejected, either block or choose any
$a'$ with $ua'\in\Good(P)$ and update the recognizer on $a'$. The committed
run remains good, and no rewrite fires on a compliant proposed trace, so
soundness and transparency hold. Conversely, the T1 limit argument and
never-alter recognizer apply to gates with substitution as well, so adding that
decision cannot enlarge the class. The oracle recognizer rejects only bad extensions, so it
never alters a compliant proposal and $U=0$ with or without substitution. The same holds for every transparent gate: if the committed
effect history is $u$ and $ua\in\Good(P)$, prefixwise transparency on the open-loop stream of a
completion $w\in P$ of $ua$ forces $\allow$ at $(u,a)$, and the decision depends only on $(u,a)$.

(ii)~Under the stated loss normalization, a compliant proposal is flagged, and so sent to
the rewrite branch, with probability $\alpha(\theta)$ and costs at most
$\ell_{\rm sub}$ there; a flagged violation is replaced by a prefix-safe action instead of
being blocked, so the decisions on violations, and hence the miss, are unchanged. Its expected conditional loss is therefore at most
$\ell_{\rm sub}\alpha(\theta)$. This is a graded loss: $U$ (Def.~\ref{def:gate}) charges every alteration, a substitution
included, as a unit loss, so substitution leaves $U$ unchanged and lowers only this graded
false-block loss. Without the prefix-safe rewrite and loss
assumptions, only the class-equivalence statement remains; substitution need not
improve an arbitrary task utility.
\end{proof}

\begin{remark}[Escalation]
Replacing a block by an oracle that always returns a prefix-safe commit action is
a special case of substitution (with the oracle's loss in place of
$\ell_{\rm sub}$). The class conclusion needs the oracle to be part of the interface $\Pi$: an
external oracle that decides $\Good(P)$ would let a trivial gate enforce every nonempty safety $P$,
beyond $\RA[\Th;\Pi]$ (the safety boundary itself still holds by the limit-point argument). The
utility conclusion is relative to that oracle model; no claim is made for a fallible or
non-prefix-safe human decision.
\end{remark}

\subsection*{A.3\quad T2: enforcement under a fallible judge}
For a one-step binary decision, let $F_1$ and $F_0$ be the laws of pre-commit
information $I$ conditional on a violating and compliant proposal. A randomized
block rule is a measurable $\varphi(I)\in[0,1]$. Its miss and false-block
probabilities are $\beta(\varphi)=1-\E_1\varphi$ and
$\alpha(\varphi)=\E_0\varphi$.

\begin{lemma}[Neyman--Pearson optimality]\label{lem:np}
Among all block rules measurable with respect to $I$ that have miss at most $\beta_0$, the minimum
false-block probability is $\ROC_I(\beta_0)$, attained by a randomized
likelihood-ratio threshold.
\end{lemma}
\begin{proof}
Blocking a violation has probability $\Pdist[\block\mid\pi=1]=\E_1[\varphi]
=1-\text{miss}$, and false-alarm is $\Pdist[\block\mid\pi=0]=\E_0[\varphi]$.
The constraint is $\E_1\varphi\ge1-\beta_0$ and the objective is
$\E_0\varphi$. This is the Neyman--Pearson
program \citep{neyman1933problem,tong2018np}; its randomized likelihood-ratio
solution has value $\ROC_I(\beta_0)$. Boundary randomization convexifies the
error region, so $\beta\mapsto\ROC_I(\beta)$ is nonincreasing and convex.
\end{proof}

\noindent\textbf{Proposition~\ref{thm:frontier}} (Exogenous information frontier)\textbf{.}
\emph{For a fixed one-step law, the exact frontier from information $I$ is
$B_I^\star(\beta)=\ROC_I(\beta)$. If $I'$ Blackwell-refines $I$, then
$\ROC_{I'}(\beta)\le\ROC_I(\beta)$ pointwise, with equality at every $\beta$ when the
likelihood ratio of $I'$ is almost surely a function of $I$.}

\begin{proof}
The equality is Lemma~\ref{lem:np}. If $I$ is a garbling of $I'$, every
$I$-measurable rule is implementable from $I'$ by applying the garbling first.
Thus the feasible error region for $I'$ contains the one for $I$, giving the
pointwise inequality. Conversely, if the likelihood ratio of $I'$ is almost surely a
function of $I$, the randomized likelihood-ratio test attaining $\ROC_{I'}$
(Lemma~\ref{lem:np}) is $I$-measurable, so $\ROC_{I'}=\ROC_I$ at every $\beta$. A strict gain
thus needs label information in $I'$ beyond $I$; conditional dependence among judges alone
does not guarantee it. When the likelihood ratio is a nondecreasing function of a score $s$, its
level sets are $s$-intervals on which $dF_1=k\,dF_0$, so blocking the upper part of the boundary
interval (randomizing only at an atom of $s$) has the same miss and false-block probability as
randomizing over the whole level set: a threshold on $s$, randomized on ties, attains $\ROC_I$.
\end{proof}

\medskip
\noindent\textbf{Corollary~\ref{cor:reduction}} (Continuation-neutral episode lift)\textbf{.}
\emph{Under the product specialization of Assumption~\ref{ass:benign}, a stationary
test with fresh per-step randomization has $R_H(\beta)=1-(1-\beta)^H$ on the
all-violation stream and $R_H(\beta)=1-(1-q\beta)^H$ for i.i.d.\
$\mathrm{Bernoulli}(q)$ labels. Within this product class,
$B_{I,H}^\star(\delta)=\ROC_I(R_H^{-1}(\delta))$.}

\begin{proof}
The product formulas follow because ``a violation is committed'' has probability
$\beta$ (or $q\beta$) independently at each step. Since the repeated test has
per-step false-block probability $\alpha$, imposing $R_H(\beta)\le\delta$ is
equivalent to $\beta\le R_H^{-1}(\delta)$; Proposition~\ref{thm:frontier} gives the optimum
at the largest admissible miss. Shared episode-level randomization is excluded
because it need not obey the product formula. Open-loop time-varying tests with per-step
misses $\beta_t$ do no better: independence gives $R_H=1-\prod_t(1-q\beta_t)$ ($q=1$ on the
all-violation stream), so $R_H\le\delta$ iff $\sum_t\gamma_t\le-\ln(1-\delta)$ with
$\gamma_t=-\ln(1-q\beta_t)$. The per-step cost $f(\gamma)=\ROC_I\big((1-e^{-\gamma})/q\big)$ is
convex, a convex nonincreasing function (Lemma~\ref{lem:np}) of a concave increasing one, so
Jensen's inequality and monotonicity of $f$ give
$\frac1H\sum_t f(\gamma_t)\ge f\big(\frac1H\sum_t\gamma_t\big)\ge f\big(-\ln(1-\delta)/H\big)
=\ROC_I(R_H^{-1}(\delta))$ for $\delta\le R_H(1)$; for larger $\delta$ allow-all is feasible. History dependence can do strictly better even without
feedback: with an uninformative score ($\ROC_I(\beta)=1-\beta$), $H\ge2$ and $0<\delta<R_H(1)$,
blocking at step one with probability $1-\delta/R_H(1)$ and repeating that decision at every later
step gives risk $\delta$ at per-step false-block probability $1-\delta/R_H(1)<1-R_H^{-1}(\delta)$,
because $R_H$ is strictly concave, so $R_H^{-1}$ is strictly convex on $[0,R_H(1)]$ with
$R_H^{-1}(0)=0$ and $R_H^{-1}(R_H(1))=1$. On the all-violation stream the per-step miss budget
$R_H^{-1}(\delta)=1-(1-\delta)^{1/H}$ is at least the union-bound budget $\delta/H$ of Theorem~\ref{thm:conformal}, because
$(1-\delta/H)^H\ge1-\delta$; the gap is what the product conditions buy.
\end{proof}

\medskip
\noindent\textbf{Theorem~\ref{thm:sequential}} (Endogenous finite-model frontier)\textbf{.}
\emph{For a specified finite model, the controlled frontier of
Definition~\ref{def:controlled} is the image of
a reachable-belief occupancy polytope; its lower boundary is convex and
piecewise linear, its Lagrange dual is solved by belief-state Bellman recursion,
and an optimum mixes at most two deterministic belief-Markov policies.}
\begin{proof}
Augment the latent state with an absorbing unsafe bit and let $r_t$ be one only
on first entry, so $\sum_t r_t$ is the episode-unsafe indicator. With finite
abstractions and finite $H$, finitely many observation histories and posteriors
$b_t=\Pdist(X_t\mid h_t)$ are reachable. The posterior is policy-independent: the factors
$\psi(D_s\mid h_s)$ in the joint law of $(X_{\le t},h_t)$ are functions of $h_t$ and cancel in Bayes'
rule. Hence $b_{t+1}$ is a fixed function of $(b_t,D_t)$ and the next observation, whose conditional
law given $(h_t,D_t)$ depends only on $(b_t,D_t)$, so $(b_t)$ is a controlled Markov chain with a
policy-independent kernel, and $\E[c_t\mid h_t,D_t]$, $\E[r_t\mid h_t,D_t]$ depend on $h_t$ only through
$b_t$. Belief is a sufficient controlled state:
the conditional action probabilities of any perfect-recall policy define a
belief-Markov policy with the same occupancy measure, hence the same $(R,C)$.

Let $\zeta_t(b,d)$ be the probability of occupying belief $b$ and choosing $d$, and $K_t(b'\mid b,d)$
the policy-independent belief transition kernel (the Bayes update after $d$ and the next observation).
The nonnegative flow constraints are
\[
\begin{aligned}
\sum_d \zeta_1(b,d)&=\Pdist[b_1=b],\\
\sum_d \zeta_{t+1}(b',d)
  &=\sum_{b,d}\zeta_t(b,d)K_t(b'\mid b,d).
\end{aligned}
\]
Conversely, disintegrating a feasible flow gives a belief-Markov policy. Both
$R$ and $C$ are linear in $\zeta$, proving the main occupancy program (if $\delta$ is below the least
achievable $R$, the program is infeasible and $F_H(\delta)=+\infty$; the claims below concern feasible
$\delta$). Its feasible
set is a bounded polytope; its two-dimensional performance image is a compact
convex polygon, with a convex piecewise-linear lower boundary.

Finite LP duality gives
\[
\begin{aligned}
F_H(\delta)
  &=\max_{\lambda\ge0}\{\textstyle\sum_{b}\Pdist[b_1=b]\,V_1^\lambda(b)-\lambda\delta\},\\
V_t^\lambda(b)
  &=\min_d\{c_t(b,d)+\lambda r_t(b,d)\\[-2pt]
  &\hspace{34mm}+\E[V_{t+1}^\lambda(b')\mid b,d]\}.
\end{aligned}
\]
With $V^\lambda_{H+1}\equiv0$, a deterministic minimizer exists for each $\lambda$. Any constrained boundary
point lies on a supporting segment of the performance polygon. Every endpoint
vertex is the image of an extreme occupancy point, and extreme finite-horizon
occupancy points are induced by deterministic belief-Markov policies. Hence
sampling one endpoint policy once per episode implements the boundary point
with support at most two.
\end{proof}
Comparing Bellman $Q$-values gives the dynamic rule in the main text:
for $c_t=w_t\mathbf 1[\block,\text{compliant}]$ with $w_t$ fixed given $b_t$, and unsafe entry only via an
allowed violation, blocking costs $w_t(1-p_t)$, allowing incurs $\lambda p_t$ (with the unsafe bit observed, so an
allowed violation is a first entry), and rearranging
$Q(\block)\le Q(\allow)$ yields
$\lambda p_t-w_t(1-p_t)\ge\Gamma_t^\lambda$.

\medskip
\noindent\textbf{Proposition~\ref{prop:feedback}} (ROC does not identify a feedback frontier)\textbf{.}
\emph{Two $H=2$ systems can have identical judge ROC and identical ungated trajectory law but different $F_H(0)$.}
\begin{proof}
Take $H=2$, a constant score that is the gate's only information about the labels, and $Y_1\sim\mathrm{Bernoulli}(q)$ with
$0<q<1$. If step one is allowed, let $Y_2=0$ in both systems, so their ungated
laws and ROCs coincide. If step one is blocked, system A has $Y_2=0$, while
system B presents $Y_2\sim\mathrm{Bernoulli}(p)$, $0<p<1$. Zero risk forces a
first-step block. A then need not block again, whereas B must. For expected
compliant-block count, set $c_t=\mathbf{1}[\text{block and compliant}]$ and let
$r_t$ indicate first unsafe entry as in Definition~\ref{def:controlled}. Then
$F_A(0)=1-q$ and
$F_B(0)=(1-q)+(1-p)$, proving non-identifiability. The constant score is not essential: if $F_0,F_1$ are not mutually singular, let $\kappa=1-\sup\{F_0(N):F_1(N)=0\}>0$; zero risk forces blocking an $F_1$-full set at every step that may violate, so $F_A(0)=(1-q)\kappa<(1-q)\kappa+(q+(1-q)\kappa)(1-p)\kappa=F_B(0)$.
\end{proof}

\medskip
\noindent\textbf{Corollary~\ref{cor:no-estimator}} (Identification impossibility)\textbf{.}
\emph{Over any model class containing systems A and B of Proposition~\ref{prop:feedback}, no estimator whose input is limited to ungated trajectories and static judge scores can uniformly identify $F_H(0)$, hence it cannot identify the full controlled frontier.}
\begin{proof}
Any estimator restricted to ungated trajectories and static scores has the
same input law in systems A and B, but the target takes the distinct values
$F_A(0)$ and $F_B(0)$. It therefore cannot identify both uniformly. Because
$F_H(0)$ is one point of the controlled frontier, neither can it identify the
full frontier.
\end{proof}

\medskip
\noindent\textbf{Proposition~\ref{prop:explore}} (Exploratory blocking restores identification)\textbf{.}
\emph{In the model of Definition~\ref{def:controlled}, let a logging gate choose $D_t$ from $h_t$
with known probabilities $\mu(d\mid h_t)\ge\varepsilon>0$ for both decisions at every reachable
history, and let the logs record $(h_t,D_t,c_t,r_t)$ at every step. Then for every randomized
perfect-recall gate $\psi$, with $W_t=\prod_{s\le t}\psi(D_s\mid h_s)/\mu(D_s\mid h_s)\le\varepsilon^{-t}$,
$R(\psi)=\E_\mu[\sum_tW_tr_t]$ and $C(\psi)=\E_\mu[\sum_tW_tc_t]$; hence the achievable $(R,C)$ set
and $F_H(\delta)$ are identified.}
\begin{proof}
Write a logged trajectory as $\omega=(h_t,D_t,c_t,r_t)_{t\le H}$, where $h_{t+1}$
appends the next observation to $(h_t,D_t)$, and let $\omega_{<t}$ be its prefix before $h_t$.
Let $\nu$ be a gate whose decision at step $t$ is drawn from $\nu(\cdot\mid h_t)$ with
randomization independent of everything else. By Definition~\ref{def:controlled}, the next
state, observation, cost, and unsafe indicator depend on the past only through $(X_t,D_t)$, and
$(c_t,r_t)$ are functions of $(X_t,D_t,X_{t+1})$. Let $G_t(\,\cdot\mid\omega_{<t},h_t,D_t)$ be
the conditional law of $(c_t,r_t,h_{t+1})$ given the logged prefix and $(h_t,D_t)$. It is
obtained by Bayes' rule from the initial law and the kernels, integrating over the posterior of
the hidden states given $(\omega_{<t},h_t,D_t)$; the factors $\nu(D_s\mid h_s)$, $s\le t$, that
appear in the joint law of the hidden states and the logged prefix are functions of the logged
prefix alone and cancel in that posterior, so $G_t$ does not depend on $\nu$. Therefore
\[
  \Pdist_\nu(\omega)=\Pdist(h_1)\prod_{t}\nu(D_t\mid h_t)\,G_t(c_t,r_t,h_{t+1}\mid\omega_{<t},h_t,D_t).
\]
Fix $t$. The pair $(c_t,r_t)$ is a function of the logged prefix $\omega_{\le t}$ ending with
$(c_t,r_t)$, whose law under $\nu$ is the product above truncated at $t$. Since
$\mu(d\mid h)\ge\varepsilon>0$ at every reachable $h$, $\Pdist_\psi(\omega_{\le t})>0$ implies
$\Pdist_\mu(\omega_{\le t})>0$, and the likelihood ratio is
$\Pdist_\psi(\omega_{\le t})/\Pdist_\mu(\omega_{\le t})=W_t$ because the common factors
$\Pdist(h_1)$ and $G_s$ cancel. Hence $\E_\psi[r_t]=\E_\mu[W_tr_t]$ and
$\E_\psi[c_t]=\E_\mu[W_tc_t]$; summing over $t$ gives both identities. Each identity holds for
every $\psi$, so the map $\psi\mapsto(R(\psi),C(\psi))$, its image, and
$F_H(\delta)=\inf\{C(\psi):R(\psi)\le\delta\}$ are functionals of the logging law. Each factor
$\psi(D_s\mid h_s)/\mu(D_s\mid h_s)$ is at most $1/\varepsilon$, so $W_t\le\varepsilon^{-t}$. By the two identities, the average of
$\sum_tW_tr_t$ (resp.\ $\sum_tW_tc_t$) over i.i.d.\ logged episodes is unbiased for $R(\psi)$ (resp.\
$C(\psi)$). The weight bound is attained, so the variance can grow exponentially in $H$: if
$\varepsilon\le1/2$, $\psi$ is deterministic, $\mu$ gives $\psi$'s choice probability exactly
$\varepsilon$ at every history, and $c_t=\mathbf 1[t=H]$, then $\sum_tW_tc_t=W_H$ equals
$\varepsilon^{-H}$ with probability $\varepsilon^{H}$ and $0$ otherwise, so its variance is
$\varepsilon^{-H}-1$.
Episode-level mixtures of such gates (as in Theorem~\ref{thm:sequential}) attain the
corresponding convex combinations of identified $(R,C)$ pairs, so they are identified as well.
\end{proof}

\medskip
\noindent\textbf{Theorem~\ref{thm:conformal}} (Distribution-free safety certificate)\textbf{.}
\emph{Under Assumption~\ref{ass:benign}, calibrating the threshold by conformal
risk control on violating examples to hold the per-step conditional miss risk
$r(\theta)=\Pdist[\allow\mid\text{violation}]$ at $\delta/H$ yields
$\Pdist[\text{unsafe episode}]\le\delta$ marginally over deployment episodes and the
calibration draw, finite-sample and distribution-free under per-step
exchangeability.}

\begin{proof}
Let $V_t$ be the event that step $t$ proposes a violation, with score $s_t$; the gate
allows it iff $s_t<\hat\theta$, where the calibrated threshold $\hat\theta$ is a function
of the calibration sample. An unsafe episode commits a violation at some step (a block commits no effect, and by
Assumption~\ref{ass:benign} committing a compliant proposal keeps the effect prefix in $\Good(P)$), so by the
\emph{union bound}, without conditioning on $\hat\theta$,
\begin{align*}
  \Pdist[\text{unsafe}]
  &\le\sum_{t\le H}\Pdist[V_t]\,\Pdist[s_t<\hat\theta\mid V_t]\\
  &\le H\max_{t\le H}\Pdist[s_t<\hat\theta\mid V_t],
\end{align*}
using $\Pdist[V_t]\le1$. Fix $t$ and condition on $V_t$. By label-conditional per-step
exchangeability (Assumption~\ref{ass:benign}), the $n$ violating calibration scores and
$s_t$ form an exchangeable $(n+1)$-tuple. The miss loss $\mathbf{1}[s<\theta]$ is bounded
and nondecreasing in $\theta$ (raising the threshold blocks fewer actions);
reparameterizing by $\lambda=-\theta$ makes it nonincreasing. Conformal risk
control \citep{angelopoulos2024conformal} at level $\delta/H$ on this tuple gives
$\Pdist[s_t<\hat\theta\mid V_t]\le\delta/H$, jointly over the calibration draw and the
step-$t$ proposal. (Equivalently, for $\hat\theta$ the $k$-th smallest calibration score
with $k=\lfloor(n+1)\delta/H\rfloor$, exchangeability of ranks gives
$\Pdist[s_t<\hat\theta\mid V_t]\le k/(n+1)$, ties included; $k=0$ means block-all.) The
exchangeability required is \emph{label-conditional} (Mondrian): since the miss is
calibrated on the \emph{violating} subpopulation, the violating calibration examples must
be exchangeable with the violating deployment steps, as required by
Assumption~\ref{ass:benign}, not merely jointly exchangeable before conditioning on the
label. Hence $\Pdist[\text{unsafe}]\le H\cdot\delta/H=\delta$, which is the stated
marginal guarantee. The argument never conditions on $\hat\theta$: without independence
between deployment and calibration, $\hat\theta$ can carry information about the
deployment law, so the conditional form $\Pdist[\text{unsafe}\mid\hat\theta]\le
H\,r(\hat\theta)$ with $\E[r(\hat\theta)]\le\delta/H$, for the population miss
$r(\theta)=\Pdist[\allow\mid\text{violation}]$, is justified when deployment
steps are drawn independently of the calibration sample (as in
Theorem~\ref{thm:adv}) but can fail otherwise. Nothing in this argument bounds the
finite-sample threshold's distance from the population ROC optimum.
\end{proof}

\medskip
The certificate is marginal over violation types: if a type makes up a fraction
$q$ of violations, its own per-step miss is guaranteed only to be at most $\delta/(Hq)$,
since $\Pdist[\allow\mid\text{violation},\text{type}]\le
\Pdist[\allow\mid\text{violation}]/q$. Calibrating separately per policy or per
argument role, as concurrent work does for tool-call fields \citep{rahman2026role},
refines this at the cost of more calibration violations per stratum.

\medskip
\noindent\refstepcounter{proposition}\textbf{Proposition~\theproposition} (Episode-level certificate, no $H$-penalty)\textbf{.}
\emph{If decisions alter neither later proposals nor their labels and scores (open loop) and a calibration set of whole
episodes is exchangeable with deployment
episodes, calibrating the threshold by conformal risk control on the monotone
episode loss $L(\text{ep})=\mathbf{1}[\text{some step allows a violation}]$ (violations as in
Assumption~\ref{ass:benign}) to level
$\delta$ gives $\Pdist[\text{unsafe episode}]\le\delta$ with no factor of $H$.}

\begin{proof}
The union-bound step above pays a factor $H$ because it controls each per-step risk
separately. Instead treat the whole episode as one calibration unit. The loss
$L(\text{ep};\theta)=\mathbf{1}[\exists t\le H:\text{a violation at }t\text{ scores
below }\theta]$ is equivalently
\[
L(\text{ep};\theta)
=\mathbf{1}\!\left\{\min_{t:Y_t=1}s_t<\theta\right\},
\]
with the minimum over an empty set defined as $+\infty$. It is bounded in
$[0,1]$ and \emph{nondecreasing} in $\theta$ (raising
$\theta$ blocks weakly \emph{fewer} actions, so it never turns a missed violation into a
caught one; this uses the open-loop hypothesis, since under feedback the gated episode's
unsafe indicator can be non-monotone in $\theta$). Reparameterizing by $\lambda=-\theta$
makes it a non-increasing admissible
loss for conformal risk control \citep{angelopoulos2024conformal}, which, on an
exchangeable episode calibration set, returns $\hat\theta$ with $\E[L(\hat\theta)]\le\delta$. Since
$\Pdist[\text{unsafe episode}\mid\hat\theta]=\E[L(\text{ep};\hat\theta)\mid\hat\theta]$,
taking the calibration expectation gives $\Pdist[\text{unsafe episode}]\le\delta$
directly. The price is that exchangeability is now required at the episode level
(episodes i.i.d.\ or exchangeable), stronger than the per-step version but free of
the $H$-fold tightening, which matters when $H$ is large.
\end{proof}

\medskip
\noindent\refstepcounter{proposition}\textbf{Proposition~\theproposition} (General region is copula-bounded)\textbf{.}
\emph{For the conservative any-flag gate and fixed marginal error rates,
Fr\'echet--Hoeffding bounds bracket the step-level miss for a specified violated
predicate set and the false-block probability on compliant actions. At a fixed
threshold vector the bounds are sharp for every $m$: each pair in the product of the
two intervals is attained by some choice of class-conditional couplings. At $m=1$
they equal the one-judge errors.}

\begin{proof}
Take the conservative gate that blocks when any judge flags. A violation with true
hazard on predicate set $S$ is committed iff no judge flags it, an event contained in
$\bigcap_{i\in S}M_i$ (every relevant judge misses) with $\Pdist[M_i]=\beta_i$, and equal
to it when $S=\{1,\ldots,m\}$, e.g., under the single predicate of
Assumption~\ref{ass:benign} (for $|S|<m$ only the upper miss bound below applies); a compliant action is altered iff
some judge false-alarms, $\bigcup_i E_i$ with $\Pdist[E_i]=\alpha_i$. For fixed
marginals the Fr\'echet--Hoeffding inequalities give
\[
  \max\big(0,{\textstyle\sum_{i\in S}}\beta_i-|S|+1\big)
   \le\Pdist\big[{\textstyle\bigcap_{i\in S}}M_i\big]\le\min_{i\in S}\beta_i,
\]
\[
  \max_i\alpha_i\le\Pdist\Big[{\textstyle\bigcup_i}E_i\Big]
   \le\min\!\Big(1,{\textstyle\sum_i}\alpha_i\Big),
\]
which hold for every coupling of the errors. Each endpoint is attained by an
explicit coupling, for every $m$. Realize the error indicators of a class as arcs
of the unit circle under a uniform point. Nested arcs $M_i=[0,\beta_i)$ and
$E_i=[0,\alpha_i)$ (the comonotone coupling) attain the intersection upper bound
$\min_{i\in S}\beta_i$ and the union lower bound $\max_i\alpha_i$. Placing the
complements $M_i^{c}$ as consecutive arcs of lengths $1-\beta_i$, wrapping around the
circle when necessary, gives
$\Pdist[\bigcap_{i\in S}M_i]=\max(0,\sum_{i\in S}\beta_i-|S|+1)$; placing the $E_i$
as consecutive arcs of lengths $\alpha_i$ gives
$\Pdist[\bigcup_iE_i]=\min(1,\sum_i\alpha_i)$. No $m$-variate countermonotone
copula is needed (none exists for $m\ge3$), because each quantity is the
probability of a single event. Misses are defined on the violation class and false
alarms on the compliant class, so the two couplings are chosen independently, and
mixing two couplings within a class attains every intermediate value. The
threshold-varying question is different: one coupling of the score vector must serve
every threshold vector at once, and for $m\ge3$ the pointwise extremes need not be
simultaneously attainable (Supplement~B.16, O1). An episode lift additionally
requires specified temporal dependence; only under Cor.~\ref{cor:reduction}'s
product conditions may these step probabilities be inserted into its product
formula. For $m=1$ both bounds equal $(\alpha_1,\beta_1)$.
\end{proof}

\medskip
\noindent\textbf{Theorem~\ref{thm:adv}} (Adversarial deployment bound)\textbf{.}
\emph{Under Assumption~\ref{ass:benign}, with conformal calibration on benign
data and a budget-$\rho$ adversary acting at deployment, on steps drawn independently of
the calibration sample and with the violating proposals generated under attack (before
perturbation) satisfying Assumption~\ref{ass:benign}, the calibrated gate satisfies
$\Pdist[\text{unsafe episode}]\le\min\{1,\delta+H\bar\eta(\rho)\}$ marginally over the calibration draw. Conditional on a fixed
threshold, an attained-margin independent product construction realizes
$1-(1-\min\{1,\beta(\theta)+\eta(\rho;\theta)\})^H$.}

\begin{proof}
\emph{Upper bound.} Calibration is performed on the benign exchangeable sample;
the representation adversary acts only at deployment, on steps drawn independently
of the calibration sample, so that conditioning on $\hat\theta$ leaves each step's
benign miss risk equal to $r(\hat\theta)$. (Without independence, the per-step
guarantee of Theorem~\ref{thm:conformal} still holds marginally, because conformal risk control is itself a marginal
guarantee; the margin argument below needs the conditional form.) By
Def.~\ref{def:adversary}, $\bar\eta(\rho)$ upper-bounds the
per-step increase in the conditional miss risk \emph{at every threshold}, so under any
budget-$\rho$ adversary the risk at the calibrated $\hat\theta$ is at most
$r(\hat\theta)+\bar\eta(\rho)$. The union bound of the certificate proof, conditioning on
$\hat\theta$, gives $\Pdist[\text{unsafe}\mid\hat\theta]\le\sum_{t\le H}(r(\hat\theta)+\bar\eta(\rho))
=H\,r(\hat\theta)+H\bar\eta(\rho)$; taking the calibration expectation and using
$\E[r(\hat\theta)]\le\delta/H$ (conformal risk control) yields the marginal
$\Pdist[\text{unsafe episode}]\le\min\{1,\delta+H\bar\eta(\rho)\}$.
\emph{Sharpness construction.} Fix $\theta$ and suppose a representation adversary
attains the defining supremum, so its conditional miss is
$x=\min\{1,\beta(\theta)+\eta(\rho;\theta)\}$. On an all-violation product stream,
choose its perturbation randomness independently across steps. The probability
that no violation is committed is $(1-x)^H$, giving exactly
$1-(1-x)^H$. If the supremum is not attained, the same risks are approached with
$\eta-\varepsilon$. This demonstrates first-order sharpness of the union bound
under a stated construction; it is not a lower bound for every adversary or
feedback process.
\emph{Non-transfer.} Conformal validity requires the deployment steps to be
exchangeable with the calibration sample; a budget-$\rho$ adversary changes the
input distribution, so the \emph{attacked} per-step miss need not satisfy the $\delta/H$
bound (the benign $\E[r(\hat\theta)]\le\delta/H$ still holds) and the bare
$\delta$ bound is void unless $\eta$ absorbs the shift.
\end{proof}

\medskip
\noindent\textbf{Proposition~\ref{prop:zero-miss}} (Distributional zero-miss price)\textbf{.}
\emph{For any attainable adversarial violation-score law $Q_\rho$, every
score-measurable gate with zero $Q_\rho$-miss has utility loss
\[U\ge1-\sup\{F_0(N):Q_\rho(N)=0\};\]
if $F_0\ll Q_\rho$, this distributional
zero-miss cost is one.}
\begin{proof}
Let $Q_\rho$ be any attainable score law
induced on violations and $F_0$ the compliant score law. Fix a score-measurable
gate with zero miss under $Q_\rho$. For a deterministic gate the blocked set
$B$ has $Q_\rho(B)=1$; for a randomized gate the same argument applies to its
block-probability function and yields the identical variational value. Its utility loss is the
compliant mass it over-blocks, $U(g)=F_0(B)$, and over all zero-miss gates
\begin{align*}
  U(g)\ &\ge\ \inf\{F_0(B):Q_\rho(B)=1\}\\
        &=\ 1-\sup\{F_0(N):Q_\rho(N)=0\},
\end{align*}
using $F_0(B)=1-F_0(B^{c})$ with $B^{c}$ ranging over the $Q_\rho$-null sets. This is an
\emph{absolute-continuity} (Lebesgue-decomposition) bound, sharper and correct where a raw support-overlap
bound is not: a zero-miss gate may leave unblocked only the part of the compliant law
\emph{singular} to $Q_\rho$. By the Lebesgue decomposition $F_0=F_0^{\mathrm{ac}}+F_0^{\mathrm{sing}}$
relative to $Q_\rho$, $U(g)\ge F_0^{\mathrm{ac}}(\mathbb{R})=1-F_0^{\mathrm{sing}}(\mathbb{R})$.
Hence if $F_0\ll Q_\rho$ (no compliant mass on any $Q_\rho$-null set) then $U(g)=1$:
block-everything is the only distributionally zero-miss score rule. This is not
the universal streamwise soundness of Def.~\ref{def:gate}. We emphasize that
$\eta(\rho;\theta)\ge1-\beta(\theta)$ is \emph{not} sufficient for domination:
it drives the adversarial miss to one at that threshold, but says nothing about
whether $F_0$ is dominated by $Q_\rho$ ($F_0$
may still place mass on $Q_\rho$-null sets, e.g., high-scoring compliant actions). What
approaches $F_0\ll Q_\rho$ is empirical: in the safe-washing runs the adversary drives $Q_\rho$
onto the low scores where the LLM judges already place compliant actions, so the $Q_\rho$-absolutely-continuous compliant mass
$1-\sup\{F_0(N):Q_\rho(N)=0\}$ is large and the zero-miss gate's over-block $U$ is near $1$; we
report the measured $U$ rather than assume total absolute continuity.
\end{proof}

\medskip
\noindent\textbf{Corollary~\ref{cor:robust}} (Robust calibration)\textbf{.}
\emph{If a uniform margin $\bar\eta(\rho)<\delta/H$ has been verified, calibrating on
benign exchangeable violations to the per-step target
$\delta/H-\bar\eta(\rho)$ yields $\Pdist[\text{unsafe episode}]\le\delta$ (marginal over the calibration
draw) under any budget-$\rho$ representation adversary at deployment, under the
conditions of Theorem~\ref{thm:adv}. At the population
optimum the frontier price is
$\ROC_I(\delta/H-\bar\eta(\rho))-\ROC_I(\delta/H)$; finite-sample calibration may pay
additional conservatism (Proposition~\ref{prop:zero-miss} prices distributional zero miss).}

\begin{proof}
Run the same split-conformal calibration but against the tightened per-step target
$\tau'=\delta/H-\bar\eta(\rho)>0$, so the benign per-step miss satisfies
$\E[r(\hat\theta)]\le\tau'$ over the calibration draw (the miss at a given calibrated
threshold can exceed $\tau'$). By
Def.~\ref{def:adversary} a budget-$\rho$ adversary raises the per-step conditional miss at
\emph{any} threshold by at most $\bar\eta(\rho)$, to at most $r(\hat\theta)+\bar\eta(\rho)$,
whose calibration expectation is at most $\tau'+\bar\eta(\rho)=\delta/H$. The
union bound of the upper part, conditional on $\hat\theta$ and then averaged over the
calibration draw, gives $\Pdist[\text{unsafe}]\le\sum_{t\le H}\delta/H=\delta$,
for \emph{any} budget-$\rho$ adversary since the per-step bound is worst-case.
At the population optimum, changing the miss target from $\delta/H$ to $\tau'$
changes false-block cost by
$\ROC_I(\tau')-\ROC_I(\delta/H)$. A threshold selected from a finite calibration
sample can lie above this envelope, so equality of its realized cost is not
claimed. If the tightened target is non-positive, this procedure falls back to
block-all; the separate domination condition $F_0\ll Q_\rho$ is what forces unit
cost for distributional zero-miss.
\end{proof}

\subsection*{A.4\quad T3: the decidability boundary}
\noindent\textbf{Proposition~\ref{thm:undec}} (Specification analysis is undecidable)\textbf{.}
\emph{Enforcement is constructive once a nonempty policy is given as a $\Th$-representable
safety automaton (T1); static analysis of an arbitrary specification is not. There is a guard
theory $\Th^{\ast}$ (equality plus two global $\mathbb N$-counters with
increment, guarded decrement, and zero-test)
for which reachability of a designated bad state in a deterministic, trim register automaton is
undecidable; consequently deciding whether a specified data-word policy is nontrivial (has a
bad prefix, i.e., ever fires), the basic sanity check a policy DSL's linter would run, is
undecidable.}

\begin{proof}
A two-counter Minsky machine $M$ has counters $c_1,c_2\in\mathbb{N}$ and a finite
program of labeled instructions of three forms,
\[
\begin{aligned}
&\ell:\ c_j\mathbin{+{+}};\ \mathrm{goto}\ \ell',\\
&\ell:\ \text{if }c_j=0\text{ goto }\ell'\text{ else }c_j\mathbin{-{-}};\ \mathrm{goto}\ \ell'',
\end{aligned}
\]
and a halt instruction. Whether $M$, started from $c_1=c_2=0$ at its initial
label, reaches its halt instruction is undecidable.

Build a deterministic register automaton $\mathcal A$ over $\Th^{\ast}$ as follows.
Its finite control is the set of $M$'s labels plus an absorbing good sink $G$;
it has two global $\mathbb N$-counters $r_1,r_2$ initialized to $0$; and the
alphabet has exactly two effect actions, the nullary $\mathsf{tick}$ and $\mathsf{safe}$ ($\noop$ leaves the state unchanged, A.1). On each
$\mathsf{tick}$, $\mathcal A$ executes the instruction at the current label: an
increment instruction updates $r_j\mathrel{+}=1$ and moves to $\ell'$; a
test-and-decrement instruction uses the $\Th^{\ast}$ zero-test guard $r_j=0$ to
branch to $\ell'$, or otherwise updates $r_j\mathrel{-}=1$ and moves to
$\ell''$. From every nonbad machine label, $\mathsf{safe}$ moves to $G$, where
both actions loop. These updates and guards are exactly the operations of
$\Th^{\ast}$, and $\mathcal A$ is deterministic because $M$ is. Designate
$M$'s halt label as the absorbing bad state $B$.

Because $M$ is deterministic and requires no input, the run of $\mathcal A$ on the
infinite $\mathsf{tick}$ stream faithfully simulates $M$: after $n$ ticks its
configuration $(\text{label},r_1,r_2)$ equals $M$'s configuration after $n$ steps.
Hence a finite $\mathsf{tick}$-prefix reaches $B$ iff $M$ halts, so bad-state
reachability in $\mathcal A$ is undecidable. The induced policy $P$ has a bad
prefix, i.e., is nontrivial (rules out at least one finite behavior), iff $B$ is
reachable iff $M$ halts. Assuming w.l.o.g.\ that $M$'s initial label is not its halt label (such machines halt trivially), $P$ is always nonempty, since taking
$\mathsf{safe}$ before $B$ yields an infinite run in $G$, and every nonbad
prefix has such a safe continuation, so $\mathcal A$ is trim. Enforcing this constructed
$\Th^{\ast}$-automaton is therefore still constructive
(T1); what is undecidable is the static question of whether an arbitrary such specification is
non-vacuous, so no DSL linter can decide it over $\Th^{\ast}$.
\end{proof}
\noindent\emph{Trimness is undecidable for $\Th^{\ast}$.} Drop the $\mathsf{safe}$ action and make the
halt label a non-bad dead end from which every action enters a fresh absorbing bad state. Every reachable non-halt
configuration keeps a non-bad $\mathsf{tick}$ successor, so the automaton is trim (\S\ref{sec:model})
iff the halt label is unreachable, i.e., iff $M$ does not halt.

\medskip
\noindent\textbf{Theorem~\ref{thm:frag}} (A decidable, expressive fragment)\textbf{.}
\emph{Over an infinite equality-data sort and a separate binary-encoded
$\mathbb N$ amount sort, for the fragment $\Th_0$ of finite, separable, disjunctive products of total equality-register
components (E), uniform independent key-local monotone-counter components (K),
and one finite monotone counter-vector component (V), bad-prefix reachability is
\PSPACE-complete, and \PSPACE-hard already for one (E) component; without (E) components, with keys over the
equality sort, and if every amount atom is a cap check, it is \textsc{np}-complete (Proposition~\ref{prop:core}). Cap
saturation is exact for (K,V), one-witness projection is exact for (K), and the
finite abstraction is searched on the fly. For trim specifications (\S\ref{sec:model}), online enforcement is polynomial per
step in the specification, current-action encoding, and $\log$ of the number of
stored keys; with saturated counters,
bit-space is linear in stored data encodings plus per-key local abstract states.}

\begin{proof}
\emph{Exact saturation.} For counter $x$, let $C_x$ be the greatest constant
against which $x$ or $x+v$ is compared (or $0$ if none) and define
$\sigma_x(n)=\min\{n,C_x+1\}$. Because every update is nonnegative, values above
$C_x$ never return below it. For each supplied amount field $v$, let $C_v$ be
the largest $C_x$ among counters whose guards or updates use that field, guess
the \emph{single} category $\tau_v(v)=\min\{v,C_v+1\}$, and derive every
counter-specific category as $\min\{\tau_v(v),C_x+1\}$. Thus all uses represent
one concrete amount, rather than independent per-counter guesses. Every allowed atomic guard
$x\bowtie c$ or $x+v\bowtie c$ has the same truth value for all concrete
valuations with the same abstraction, and abstract addition
$\sigma_x(\sigma_x(x)+\min\{\tau_v(v),C_x+1\})$ exactly matches concrete addition followed
by $\sigma_x$. Boolean combinations and deterministic transitions therefore
preserve this equivalence. A (V) component with $k$ counters has at most
$|Q|\prod_x(C_x+2)$ abstract configurations. One transition may update several
counters, but each atomic guard still compares one counter (or that counter
plus the supplied amount) with a constant.

\emph{One-witness projection.} A (K) component has a local finite state and
counter vector for each key, initialized identically; an action naming $k$ reads
and updates only the component at $k$, and the component rejects iff some local
copy reaches its absorbing bad state. If a concrete run rejects, project the
subsequence of actions naming the first rejecting key: locality makes this
exactly a one-key rejecting run (actions naming no key update no local copy and are dropped).
Conversely, rename the key in any one-key
rejecting run to a fresh concrete key to obtain a global rejecting run. Thus
bad-prefix reachability is equivalent to reachability in one saturated local
copy, not merely over-approximated.

\emph{Equality and products.} Because the equality-data sort is infinite, an
(E) component has the standard finite equality-type abstraction for its
registers and current data; its nonemptiness problem is
decidable \citep{kaminski1994finite}, and in polynomial space by the on-the-fly search below, since each equality type has a polynomial-size encoding. Nonemptiness of register automata is itself \PSPACE-complete \citep{demri2009ltl}, so the (E) components add no complexity beyond that class. Every component is
total over the common alphabet: guards partition the input valuations and any
missing case is completed to its absorbing bad state. Separability excludes
cross-component guards and assumptions. Hence a word rejecting in one
component is executable by every other component and rejects the disjunctive
product; conversely every product rejection has such a witness component.
One may therefore nondeterministically select a component and search its
implicit abstract graph. A node stores finite control, equality type, and each
saturated counter in binary; a successor guesses the finite action/equality
type and one capped category per supplied amount field, shared across every
use of that field. The graph can be exponential
because constants are binary, but reachability uses polynomial space by an
on-the-fly counter bounded by the number of abstract nodes. By
\textsc{npspace}=\textsc{pspace}, nontriviality is in \textsc{pspace}.

\emph{Hardness.} Let $L\in\PSPACE$ be decided by a deterministic Turing machine $M$ with states
$Q$ and tape alphabet $\Gamma$ in space $p(n)$; we assume $M$ never leaves cells $1,\ldots,p(n)$.
Given $x$, build in polynomial time one total deterministic (E) component $\mathcal A_x$ over a
single unary action type, with registers $z_\gamma$ ($\gamma\in\Gamma$) and $y_1,\ldots,y_{p(|x|)}$; their number is part of the input and, as in \citet{demri2009ltl} but unlike the pairwise-distinct registers of \citet{kaminski1994finite}, two registers may hold equal data.
The first $|\Gamma|$ actions store pairwise distinct data in the $z_\gamma$, and the next $p(|x|)$
actions store in $y_i$ a datum equal to $z_{\gamma_i}$, where $\gamma_1\cdots\gamma_{p(|x|)}$ is $x$
padded with blanks. The control then holds $M$'s state $q$ and head position $h$. Let $\gamma$ be
the symbol with $y_h=z_\gamma$ (an equality guard) and $(q',\gamma',\pm1)$ the move of $M$ on
$(q,\gamma)$. An action whose datum equals $z_{\gamma'}$ stores it in $y_h$ and moves the control to
$(q',h\pm1)$. Every other datum, i.e., a repeated value during setup or a wrong symbol, leads to an
absorbing good sink, as do $M$'s rejecting states, while $M$'s accepting state is the bad state.
Guards thus partition every input, so $\mathcal A_x$ is total, deterministic given the datum's
equality type, trim, and nonempty. It has $|\Gamma|+p(|x|)$ registers and
$O(|\Gamma|+p(|x|)+|Q|\,p(|x|))$ control states, and its bad state is reachable iff $M$ accepts $x$.
Bad-prefix reachability is therefore \PSPACE-hard already for a single (E) component (cf.\ the
\PSPACE-completeness of register-automaton nonemptiness \citep{demri2009ltl}), and $\Th_0$
nontriviality is \PSPACE-complete.

\emph{Runtime cost.} The T1 gate stores $\sigma_x(x)$, not an unbounded concrete
counter, and maintains a balanced dictionary from each encoded key to its local
finite state and saturated counter vector. On an action of bit length $L$, it
reads each supplied amount once, computes its shared capped category, looks up
the touched key using $O(\log n)$ encoded-key comparisons, evaluates the guard
partition, and performs saturated updates. Time is polynomial in $L$ and the
specification, with the stated $\log n$ dictionary factor. If a saturated local
state needs $s$ bits, total (K) storage is
$O(\sum_{k\ {\rm seen}}|k|+ns)$ bits, plus the bounded E-register encodings and
fixed V abstract state. This is the claimed linear dependence on stored
encodings, not a unit-cost claim for arbitrary-precision values or hashing.

\emph{Expressiveness.} We exhibit each policy as a deterministic $\Th_0$
automaton with absorbing bad state.
\emph{Authenticate-before-access}: a control flag $\mathit{auth}$ and an
equality register $\mathit{prin}$ storing the authenticated principal; the
$\mathsf{authenticate}(p)$ action sets $\mathit{auth}$ and $\mathit{prin}:=p$; an
$\mathsf{access}$ by requester $q$ takes the bad transition unless $\mathit{auth}$
holds and $q=\mathit{prin}$ (an equality guard).
\emph{Per-recipient spend cap}: one monotone accumulator $s_k$ per recipient
key $k$ (stored in the balanced key dictionary); $\mathsf{pay}(k,v)$ with a nonnegative integer amount
$v$ takes the bad transition if $s_k+v>c$ and otherwise updates $s_k\mathrel{+}=v$, an
add-nonnegative-amount update and a compare-to-constant guard, both in $\Th_0$. Static
decidability applies the exact-saturation abstraction (exact up to the cap $c$, then absorbing
at $c{+}1$) to the single representative key, since distinct keys evolve independently, so ``some
key exceeds its cap'' is reachable iff it is reachable for one key. P2 uses the
same monitor with composite key $(k,w)$ when a trusted clock adapter attaches
the day bucket $w$.
\emph{Rate limit $r$ per supplied bucket}: key a monotone accumulator by
the action's composite $(k,w)$ field; its $(r{+}1)$-th action takes the bad
transition. The automaton enforces this policy for every supplied bucket
sequence and assumes no ordering between bucket identifiers. Interpreting
$w$ as a wall-clock window additionally requires a trusted adapter to attach
$w=\lfloor t/W\rfloor$; that adapter is outside $\Th_0$, rather than an
unstated trace assumption. No decrement or reset is needed because a new
bucket is a fresh key.
\emph{No-send-after-revoke}: a boolean control flag set by $\mathsf{revoke}$;
$\mathsf{send}$ takes the bad transition once the flag is set.
\emph{Fixed-vector global quotas}: a (V) component keeps one counter per
resource class, may test a proposed nonnegative amount against the corresponding
remaining cap, and adds the tested amount to each tested coordinate (possibly also to other coordinates). This permits coupling
within the declared finite vector but never allocates or compares counters for
unboundedly many data keys.
\end{proof}

\begin{proposition}[Policy inclusion and redundancy]\label{prop:incl}
Let $\mathcal A,\mathcal B$ be trim $\Th_0$ specifications of policies $P_{\mathcal A},P_{\mathcal B}$ over a common alphabet, where $\mathcal A$ has no (K) component
and registers, keys and equality guards range over the equality sort (amount fields are read only by counter guards and updates).
Deciding $P_{\mathcal A}\subseteq P_{\mathcal B}$, i.e., whether rule $\mathcal B$ is redundant given $\mathcal A$ ($P_{\mathcal A}\cap P_{\mathcal B}=P_{\mathcal A}$), is \PSPACE-complete,
and \PSPACE-hard already when $\mathcal A$ allows every action and $\mathcal B$ is one (E) component. Inclusion for $\mathcal A$ with a (K) component is not covered.
\end{proposition}
\begin{proof}
\emph{Reduction.} For any $P$, $u\in\Good(P)$ iff $uy\in P$ for some $y$; as $P_{\mathcal B}$ is safety, $P_{\mathcal A}\subseteq P_{\mathcal B}$ iff
$\Good(P_{\mathcal A})\subseteq\Good(P_{\mathcal B})$ iff $P_{\mathcal A}\cap P_{\mathcal B}=P_{\mathcal A}$, i.e., adding rule $\mathcal B$ changes no decision of the T1 gate.
By trimness, $\Good(P_{\mathcal A})$ ($\Good(P_{\mathcal B})$) is the set of words on which no component of $\mathcal A$ ($\mathcal B$) enters a bad state, and
$\mathcal B$ rejects iff some component $\mathcal B_j$ does. So inclusion fails iff, for some $j$, some word keeps every component of $\mathcal A$ non-bad
while $\mathcal B_j$ enters a bad state. By the typing hypothesis, equality types and amount categories never constrain the same value.
\emph{Product search.} Fix $j$. If $\mathcal B_j$ is a (K) component, add an analysis-only register tuple $r^\ast$, never assigned, whose initial
equality type relative to the initial register contents is guessed, and track only the local copy of key $r^\ast$: it starts in the initial local
state and is updated exactly by the actions whose key equals $r^\ast$. A node stores the control states of all components of $\mathcal A$ and of
$\mathcal B_j$ (or of the $r^\ast$ copy), the joint equality type of all their registers and $r^\ast$, and every counter, saturated as in the proof of
Theorem~\ref{thm:frag}; each amount field $v$ gets one shared category $\min\{v,C_v+1\}$, with $C_v$ the largest $C_x$ over counters $x$ of
$\mathcal A$ or $\mathcal B_j$ that use $v$. A successor guesses an action type, the equality type of its data relative to the registers and $r^\ast$, and
the amount categories; it is discarded if a component of $\mathcal A$ enters a bad state, and accepting if $\mathcal B_j$ (its $r^\ast$ copy) does.
\emph{Exactness.} Saturation preserves the truth value of every atom $x\bowtie c$ and $x+v\bowtie c$ of either specification and commutes with
the nonnegative updates (proof of Theorem~\ref{thm:frag}). Over the infinite equality sort, fresh values realize every consistent type sequence,
and all register, key and $r^\ast$ tests are equalities. If a word keeps $\mathcal A$ non-bad while the copy of $\mathcal B_j$ at key $k^\ast$ rejects,
then $r^\ast:=k^\ast$, with its actual initial type, maps its run to an accepting path; conversely, an accepting path realizes to a word on which
every component of $\mathcal A$ evolves as tracked and the copy at the realized $r^\ast$ rejects. No action is deleted: actions naming other keys stay
as actions with key $\neq r^\ast$ and update $\mathcal A$ exactly, so $\mathcal A$'s guards need not be monotone. Nodes have polynomial size, so an
on-the-fly search over all $j$, with a step counter bounded by the number of nodes, puts non-inclusion in \textsc{npspace}; by Savitch's theorem
and closure under complement, inclusion is in \PSPACE{}.
\emph{Hardness.} Let $\mathcal A$ be the one-state, register-free (E) component that allows every action; it is trim and $P_{\mathcal A}$ contains
every trace, so $P_{\mathcal A}\subseteq P_{\mathcal B}$ iff $\mathcal B$ has no reachable bad state. Bad-prefix reachability is \PSPACE-hard already for
one trim (E) component (proof of Theorem~\ref{thm:frag}), and \PSPACE{} is closed under complement, so inclusion is \PSPACE-hard.
In \S\ref{sec:q5}, $\mathcal A$ is the session limit, a (V) component, and $\mathcal B$ is the per-recipient cap, a (K) component; that check is an instance.
\end{proof}

\begin{proposition}[Caps, quotas, and revocation: an \textsc{np}-complete core]\label{prop:core}
Call an amount atom $x+v\bowtie c$ in the guard of a transition case (a cell of the guard partition for one control state and action
type, with its target and updates) a \emph{cap check} if that case's update of $x$ contains the summand $v$ (so $x$ grows by at
least $v$) or the case enters the absorbing bad state; completed missing cases enter the bad state, so their atoms are cap checks. Consider
$\Th_0$ specifications without (E) components whose keys range over the equality sort (amount fields are read only by counter
guards and updates, as in Proposition~\ref{prop:incl}) and in which every amount atom is a cap check. For them, bad-prefix
reachability (nontriviality) is \textsc{np}-complete; it is \textsc{np}-hard already for one (V) component with unit increments and
the constant $1$. When neither
$\mathcal A$ nor $\mathcal B$ has (E) components, $\mathcal A$ has no (K) component, and every amount atom is a cap check, deciding
$P_{\mathcal A}\subseteq P_{\mathcal B}$ for trim $\mathcal A,\mathcal B$ (Proposition~\ref{prop:incl}) is co\textsc{np}-complete. All these rules of Theorem~\ref{thm:frag} are of this form: per-recipient caps, bucket quotas, rate limits,
no-send-after-revoke (its flag a (V) counter $\varrho$ with $\mathsf{revoke}:\varrho\mathrel{+}=1$ and $\mathsf{send}$ bad if $\varrho\ge1$, or
a control bit of the (V) component), and fixed-vector quotas in which every case that tests a coordinate against a supplied amount
$v$ either enters the bad state or adds $v$ to that coordinate. So are P2, P3, and the main-paper \S\ref{sec:q5} redundancy check. Several global rules (a session limit, a revocation
flag, fixed-vector quotas) form one (V) component whose control is the product of theirs; the bounds are in the size of that
merged component, set-once flags encoded as counters tested against $1$ keep it polynomial, and nontriviality alone needs no
merging, since a disjunctive product is checked component by component. The cap-check hypothesis cannot be dropped: in one
register-free (V) component, cells $x_h+v=K\wedge y_h+v=K-b$ test $x_h-y_h=b$, so each counter pair stores one rewritable bit, and
with $K=2^{O(p(n))}$ in binary such a component simulates a polynomial-space Turing machine; without (E) components but with
arbitrary amount atoms, bad-prefix reachability therefore stays \PSPACE-complete.
\end{proposition}
\begin{proof}
\emph{Reduction to one counter component.} By the product argument and the one-witness projection in the proof of
Theorem~\ref{thm:frag}, it suffices to decide whether the bad state of one component without registers and keys (a (V) component,
or one local copy of a (K) component) is reachable. Such a component has finite control $Q$ and monotone counters $X$. For a counter
$x$, let $c^x_1<\cdots<c^x_{m_x}$ be the constants of its atoms $x\bowtie c$ and $x+v\bowtie c$, and call the $m_x$ singletons
$\{c^x_i\}$ and the nonempty gaps between consecutive constants (below the first, above the last), at most $2m_x+1$ sets in all,
the \emph{regions} of $x$. An atom
$x\bowtie c$ has one truth value on each region of $x$, and an atom $x+v\bowtie c$ has one truth value on each region containing
$x+v$. A \emph{region vector} $r$ assigns each counter a region.

\emph{Phases.} Updates are nonnegative, so along a run every counter visits its regions in increasing order, and the region vector
changes at most $R=\sum_x 2m_x$ times; $R$ is polynomial in the specification although constants are binary. A run that reaches the
bad state therefore splits into at most $R+1$ \emph{phases} of constant region vector, joined by at most $R$ \emph{boundary} steps (each
changes some region) and followed by the final step into the bad state. Consider a step inside a phase with vector $r$ that takes
case $t$. It does not enter the bad state, its pre- and post-values lie in $r$, and for every amount atom $x+v\bowtie c$ of its guard,
a cap check, $t$ adds at least $v$ to $x$, so $x+v$ lies between the pre- and post-value of $x$ and hence in $r(x)$. Every atom of
$t$'s guard thus has its value under $r$, so $t\in E_r$, the set of non-bad cases whose guard is true when each atom $x\bowtie c$ is
evaluated at $r(x)$ and each $x+v\bowtie c$ with $x+v\in r(x)$; $E_r$ is computable in polynomial time from $r$. Phase $i$ ends in the
source state of boundary case $i$, whose target starts phase $i+1$; the first phase starts in the initial state; empty phases are
allowed; and cases that cannot stay in $r$ are forced to occur zero times by the end-value bounds.
Conversely, take any path of cases in $E_r$ from control state $q$ to $q'$, with any nonnegative amounts, started from counter values in
$r$: if the values after the path are still in $r$, then, as counters are monotone and regions are intervals, every intermediate
value is in $r$, every guard on the path holds, and the path is executable.

\emph{One existential Presburger formula.} A phase is thus described by its start and end control states, the Parikh vector
$\varphi\in\mathbb N^{E_r}$ of a path between them, and, for each case $t$ and amount field $f$ of $t$, the total $A_{t,f}\ge 0$ supplied
to $f$ over $t$'s occurrences ($A_{t,f}=0$ if $\varphi(t)=0$, and any value otherwise, since one occurrence may carry it all). End
values are linear: $x^{\rm end}=x^{\rm start}+\sum_t\varphi(t)\,k_{t,x}+\sum_{t,f:\ t\text{ adds }f\text{ to }x}A_{t,f}$, with $k_{t,x}$
the fixed increment. That $\varphi$ is the Parikh vector of a path from $q$ to $q'$ in the graph $E_r$ (flow conservation and
connectivity of its support) has a polynomial-size existential Presburger definition \citep{verma2005parikh}. Guess the number of
phases, their region vectors, start and end control states, the boundary cases, and the final case into the bad state; then require
the phase constraints, membership of each phase's start and end values in its regions (interval bounds), the boundary and final
updates with explicit amount variables, the boundary and final guards as Boolean combinations of linear atoms in those explicit values,
and initial values $0$. By the two directions above, this polynomial-size existential Presburger formula is satisfiable iff the bad
state is reachable, and satisfiability of such formulas is in \textsc{np} \citep{scarpellini1984presburger}. (Equivalently, with the guesses encoded by Boolean selector
variables, reachability reduces to a single quantifier-free linear-integer-arithmetic query of the kind SMT solvers accept.) For inclusion, apply the same encoding to the product
search in the proof of Proposition~\ref{prop:incl}. As $\Th_0$ has at most one (V) component and $\mathcal A$ has no (E) or (K) component,
$\mathcal A$ is a single (V) component (or empty), so for each $j$ the product is one register-free counter component with control
$Q_{\mathcal A}\times Q_{\mathcal B_j}$ and counters $X_{\mathcal A}\uplus X_{\mathcal B_j}$. If $\mathcal B_j$ is a (K) rule tracked at $r^\ast$, split each
keyed action type into a copy naming $r^\ast$, where both components take a case, and a copy not naming it (or a keyless type), where
$\mathcal A$ takes a case and the tracked copy is unchanged. Product cases are pairs of cases, so there are polynomially many, and amounts
are shared within a step. Delete every product case in which $\mathcal A$ enters its bad state, including at the final step; the
accepting cases are those in which the tracked copy enters its bad state. Every amount atom of a remaining non-accepting case comes from a
case of $\mathcal A$ or $\mathcal B_j$ that does not enter its bad state, so it is a cap check; hence the phase argument applies,
non-inclusion is in \textsc{np}, and inclusion is in co\textsc{np}. The restriction to one (V) component matters: if $\mathcal A$ could be a product of several
register-free control components, non-inclusion would encode nonemptiness of an intersection of finite automata and be \PSPACE-hard.

\emph{Hardness.} Reduce CNF satisfiability with variables $y_1,\ldots,y_n$ and clauses $C_1,\ldots,C_m$ to one (V) component with
counters $x^T_i,x^F_i$, unit increments, and atoms comparing a counter with $1$. In control state $S$, actions $\mathsf{setT}_i$ and
$\mathsf{setF}_i$ increment $x^T_i$ and $x^F_i$, respectively, and loop, and $\mathsf{go}$ moves to $K_1$. In $K_j$, the action $\mathsf{pick}_{j,\lambda}$ for
a literal $\lambda$ of $C_j$ moves to $K_{j+1}$ if the counter of $\lambda$ ($x^T_i$ for $y_i$, $x^F_i$ for $\neg y_i$) is at least $1$, and to
an absorbing good sink $G$ otherwise. In $L_i$ (with $L_1=K_{m+1}$), $\mathsf{cons}_i$ moves to $G$ if $x^T_i\ge1\wedge x^F_i\ge1$ and to
$L_{i+1}$ otherwise, and in $L_{n+1}$ the action $\mathsf{fire}$ enters the bad state. Every other pair of control state and action moves
to $G$, which loops on every action, so the guards partition every input and the component is total, deterministic, and trim. A word
reaches the bad state iff the counters set before $\mathsf{go}$ define a consistent partial assignment that satisfies every clause, i.e.,
iff the formula is satisfiable; the construction is polynomial (it also works as one (K) component in which every action type carries the key field; a witness names a single key).
For inclusion, let $\mathcal A$ be the one-state (V) component with an empty counter vector whose single case loops on every
action; it is total, deterministic, trim, and has no (E) or (K) component. As $P_{\mathcal A}$ contains every trace and $\mathcal B$ is trim,
$P_{\mathcal A}\subseteq P_{\mathcal B}$ iff $\mathcal B$ has no reachable bad state, so inclusion is co\textsc{np}-hard.
\end{proof}


 
